\section{Related Work}

The growing computational demands of training large foundation models have driven significant research into performance modeling and optimization of DLT workloads. 

\vheading{Kernel runtime prediction} Habitat~\cite{geoffrey2021habitat} extrapolates single-GPU measurements to predict cross-device performance. More recent approaches such as NeuSight \cite{neusight}, Omniwise \cite{omniwise} rely on a mix of profiling data and architectural details of the accelerator to predict kernel runtimes more accurately. ASTRA-sim~\cite{rashidi2020astra} focuses specifically on network topology and collective communication modeling. These are complementary to \sysname and can be plugged in as needed, enabling end-to-end estimates on a wide range of workloads.

\vheading{Analytical Performance Models} Analytical models predict DLT performance through mathematical formulations of system behavior. AMPed~\cite{moolchandani2023amped} and Calculon~\cite{isaev2023calculon} propose specialized models for LLMs but support only limited parallelization strategies and require explicit modeling of new optimizations. Other work focuses on specific architectures like CNNs~\cite{yan2015performance,qi2017paleo,gianniti2018performance}. While these techniques can provide quick estimates, their applicability is limited to specific models and configurations.

\vheading{Domain-Specific Simulators} Simulation-based approaches aim to capture detailed system behavior through explicit modeling. Proteus~\cite{duan2023proteus} introduces a strategy tree abstraction for modeling parallelization patterns but requires translation into a custom specification language. DistIR~\cite{santhanam2021distir} proposes an intermediate representation for distributed computations but struggles with complex parallelization strategies. Daydream~\cite{zhu2020daydream} captures dependency graphs from execution traces, but requires GPU access and manual optimization modeling. vTrain~\cite{bang2023vtrain} uses CUPTI profiling to measure kernel runtimes but faces challenges modeling communication patterns in complex parallelization strategies. 



  

 
